%% ma: L10-B13-0003
%% lop: 10
%% chuong: 5
%% bai: 13
%% dang: 10.13.2
%% loai: TL
%% muc: TH
%% dap_an: "-"
%% nguon: "chuyen-de-cac-so-dac-trung-cua-mau-so-lieu-khong-ghep-nhom-toan-10-chuong-trinh-moi.pdf (GV. Phan Nhật Linh) (P02), Bài 02 Dạng 1 Bài tập tự luận 3"
%% kien_thuc: "Lớp 10 Bài 13 (số trung bình)"
%% trang_thai: chinh-thuc
%% ghi_chu: "Tự duyệt theo yêu cầu giáo viên 22/10/2026."
\begin{ex}
Điểm thi môn Ngữ văn của lớp 10A được cho trong bảng sau:
\begin{center}
\begin{tabular}{|c|c|c|c|c|c|c|}\hline
Điểm thi & $5$ & $6$ & $7$ & $8$ & $9$ & $10$\\\hline
Tần số & $5$ & $7$ & $12$ & $14$ & $3$ & $4$\\\hline
\end{tabular}
\end{center}
Tính điểm trung bình môn Ngữ văn của lớp 10A (làm tròn đến hàng phần mười).
\loigiai{$n=45$; $\overline x=\dfrac{5\cdot5+7\cdot6+12\cdot7+14\cdot8+3\cdot9+4\cdot10}{45}=\dfrac{330}{45}\approx7{,}3$.}
\end{ex}
